\scanpage{46}\mspage{52}
de localisation Der 2, il suffit de prouver que $i^*(u)$ et $j^*(u)$ sont iso. Pour $i^*(u)$ cela résulte du fait que $i^*$ soit foncteur de localisation, et pour $j^*$ du \deleted{\unsure{l'h}} fait que $j^*(F)$ et $j^*(i_*i^*(F))$ sont des objets finaux dans $\mathbb D(U)$.

% slanted
\margin{\deleted{Axiome de comparaison par cofinalité cohomologique}}
\deleted{\uline{Dér 6} (g) : Soit $f:X\to Y$ \deleted{un morphisme} \inserted{dans Diag, tel} que $f$ soit une \unsure{hot-|hoto-}équivalence co-stricte (i.e. $f^\circ$ \unsure{hot|hoto} équiv. stricte). Alors $\forall F\in\mathbb D(Y)$, l'hom can
\[
 f^*:\Gamma_{\mathbb D}(Y,F)\to\Gamma_{\mathbb D}(X,f^*F)\tag*{(*)}
\]
est un iso.

\uline{NB} On a un énoncé plus fort si $f$ \inserted{elle-\unsure{=|même}} est une hot équivalence stricte (pas \unsure{néces.|néc.} costricte), cf Cor. 2 p. 47.

\deleted{\uline{\unsure{Question}}} On peut se demander de caractériser les
\[
 f:X\to Y
\]
de Diag, tels que $\forall$ dérivateur exact à gauche $\mathbb D$ (ou ne serait-ce que pour \deleted{\unsure{$\mathrm{HOTAb}$}} $\mathrm{Der}(\mathrm{Ab})$ et pour $\mathrm{HOT}$), et tt $F\in\mathbb D(Y)$, l'hom $(*)$ ci-dessus soit iso. \uline{Il n'est pas vrai} que $f$ soit une hoto-équivalence costricte. (soit \deleted{\unsure{coho}} cohomologiquement cofinal), car sinon une hot-équiv. costricte serait stricte, et l'inv. serait vrai par dualité, mais ce n'est pas vrai.}
% archive sheet number: 45
